\chapter{Deep constrained monotonic networks --- Runje et al.}

This chapter extends the Runje--Shankaranarayana (2023) constrained-monotonic architecture (their
absolute-mode dense layer, split into convex and concave neuron blocks) with \emph{skip
connections}: gated residual blocks that stack the constrained-monotone dense layer to arbitrary
depth while staying trainable, and remain sound (monotone) and universal (retain the depth-4 UAP)
at every depth. The result (forthcoming) is presented in three stages: the constrained monotone
dense layer itself (\S\ref{sec:rs-dense}), the residual/deep-stack combinator and its soundness at
any depth (\S\ref{sec:skip}), and the resulting universal approximation for deep monotone networks
(\S\ref{sec:deep-uap}). The embedding of non-monotone features into a partially monotone network
--- the original framing of this chapter --- is presented as a secondary application
(\S\ref{sec:partial-mono}), extended to use the deep-residual core.

\section{Constrained monotone dense layers}\label{sec:rs-dense}

\begin{lemma}[R--S dense layer is monotone]\label{lem:rs-dense}
  \lean{UniversalApproximation.Runje.rsDense_monotone}\leanok
  \uses{def:mononet}
  The Runje--Shankaranarayana absolute-mode dense layer --- nonnegative effective weights $|W|$,
  a first block of \emph{convex} neurons through a monotone base activation $\rho$ (e.g.\ ELU,
  softplus) and a second block of \emph{concave} neurons through its point reflection
  $\mathrm{reflect}\,\rho$, appended --- is monotone.
\end{lemma}

\section{Skip connections and deep monotone networks}\label{sec:skip}

\begin{definition}[Residual block]\label{def:residual}
  \lean{UniversalApproximation.Runje.residual}\leanok
  A residual block combines a skip connection with a sublayer through nonnegative scalar gates:
  $x \mapsto g_\alpha\,\mathrm{skip}(x) + g_\beta\,F(x)$, with $g_\alpha,g_\beta \ge 0$.
\end{definition}

\begin{theorem}[Deep soundness: a residual stack of any depth is monotone]\label{thm:resnet-mono}
  \lean{UniversalApproximation.Runje.ResNet.monotone_toFun}\leanok
  \uses{def:residual}
  A \texttt{ResNet} --- a composition of residual blocks, of any depth --- in which every block has
  nonnegative gates and monotone skip/sublayer maps, is itself monotone.
\end{theorem}

\section{Deep monotone universal approximation}\label{sec:deep-uap}

\begin{theorem}[Deep monotone UAP retains UAP]\label{thm:deep-mono-uap}
  \lean{UniversalApproximation.Runje.deep_monotone_approximation}\leanok
  \uses{thm:resnet-mono, thm:mono-approx}
  Every continuous, coordinatewise-monotone $f$ on the unit cube is uniformly
  $\varepsilon$-approximated by a monotone \texttt{DeepMonoNet} (a residual body of any depth plus a
  nonnegative-weight affine read-out). The witness embeds the depth-4 \texttt{MonoNet} of
  \cref{thm:mono-approx} as a single residual block, so no depth beyond 4 is needed to retain UAP.
\end{theorem}

\section{Partial monotonicity (secondary)}\label{sec:partial-mono}

This section presents the chapter's original application: embedding non-monotone features into a
partially monotone network, extended so the monotone core may be an arbitrary-depth deep-residual
network (\cref{thm:deep-part-uap}) rather than only the shallow depth-4 core.

\subsection{Embedding non-monotone features}

A \emph{partially monotone} target $f(u,x)$ is continuous and, for each fixed non-monotone block
$u$, monotone in the monotone block $x$; its dependence on $u$ is otherwise arbitrary. The
architecture embeds $u$ through an \emph{unconstrained} single-hidden-layer network, clamps the
embedding into $[0,1]$, concatenates it with $x$, and feeds the result to a monotone network.

\begin{definition}[Unconstrained embedding span]\label{def:genspanpi}
  \lean{UniversalApproximation.Runje.genSpanPi}\leanok
  $\mathrm{genSpanPi}\,\sigma\,d_f$ is the linear span, inside all functions
  $(\mathrm{Fin}\,d_f\to\mathbb{R})\to\mathbb{R}$, of the ridge units
  $x\mapsto\sigma\!\left(\sum_c w_c x_c + b\right)$ --- the Leshno single-hidden-layer family
  written with the explicit dot product, so no inner-product instance on the $\Pi$-type is needed.
\end{definition}

\begin{definition}[Partial-monotone network]\label{def:partmononet}
  \lean{UniversalApproximation.Runje.PartMonoNet}\leanok
  \uses{def:mononet, def:genspanpi}
  A \texttt{PartMonoNet} bundles an embedding width, an unconstrained embedding
  $\mathrm{emb}\colon(\mathrm{Fin}\,d_f\to\mathbb{R})\to(\mathrm{Fin}\,N\to\mathbb{R})$, and a
  monotone network $\mathrm{mono}\colon\texttt{MonoNet}\,(N+d_m)$. Its denotation is
  $\mathrm{toFun}(u,x)=\mathrm{mono}\big(\mathrm{append}\,(\mathrm{clamp}_{[0,1]}\circ
  \mathrm{emb}\,u)\;x\big)$, where $\mathrm{clamp}_{[0,1]}$ is a fixed bounded output activation on
  the embedding.
\end{definition}

\subsection{Soundness and universal approximation}

\begin{theorem}[Soundness: monotone in $x$]\label{thm:runje-sound}
  \lean{UniversalApproximation.Runje.PartMonoNet.monotone_snd}\leanok
  \uses{def:partmononet, def:mononet}
  If the core \texttt{MonoNet} is monotone, then for every fixed $u$ the denotation
  $x\mapsto\mathrm{toFun}(u,x)$ is monotone in the (coordinatewise) $\Pi$-order.
\end{theorem}
\begin{verbatim}
theorem PartMonoNet.monotone_snd {df dm} (P : PartMonoNet df dm)
    (h : P.mono.IsMonotone) (u : Fin df → ℝ) : Monotone (P.toFun u)
\end{verbatim}
\begin{proof}\leanok
  \uses{def:partmononet}
  The clamped embedding $\mathrm{clamp}_{[0,1]}\circ\mathrm{emb}\,u$ is a fixed vector, so
  appending it and varying $x$ is monotone in the $\Pi$-order; compose with the monotone
  denotation of \texttt{MonoNet}. No continuity of $\mathrm{emb}$ is required.
\end{proof}

\begin{lemma}[Leshno embedding bridge]\label{lem:leshno-bridge}
  \lean{UniversalApproximation.Runje.leshno_bridge}\leanok
  \uses{def:genspanpi, thm:leshno-dense}
  If $\sigma$ densely approximates (Leshno), then for every compact $K\subseteq\mathrm{Fin}\,d_f\to
  \mathbb{R}$ and every function continuous on $K$, there is a total $g\in\mathrm{genSpanPi}\,
  \sigma\,d_f$ approximating it to any accuracy on $K$. Proved by transporting Leshno's density
  across the $\mathrm{EuclideanSpace}\,\mathbb{R}\,(\mathrm{Fin}\,d_f)\cong(\mathrm{Fin}\,d_f\to
  \mathbb{R})$ isometry (the Euclidean inner product equals the dot product).
\end{lemma}
\begin{proof}\leanok
  \uses{thm:leshno-dense}
  Pull the target back through the isometry to a continuous function on the (compact) preimage,
  apply Leshno density there to get a spanning element, and push the finite ridge combination back
  to a total $\Pi$-side function; the per-generator identity is exactly the inner-product / dot
  product equality.
\end{proof}

\begin{theorem}[Partial-monotone universal approximation]\label{thm:runje-uap}
  \lean{UniversalApproximation.Runje.partial_monotone_approximation}\leanok
  \uses{def:partmononet, def:genspanpi, lem:leshno-bridge, thm:mono-approx}
  For any Leshno-admissible activation $\sigma$ (class $M$, not a.e.\ polynomial) and any jointly
  continuous $f$ that is coordinatewise monotone in $x$ on the unit cubes, and any
  $\varepsilon>0$, there is a \texttt{PartMonoNet} whose core is monotone, whose every embedding
  coordinate lies in $\mathrm{genSpanPi}\,\sigma\,d_f$, and which approximates $f$ within
  $\varepsilon$ uniformly on the cube.
\end{theorem}
\begin{verbatim}
theorem partial_monotone_approximation {df dm : ℕ}
    (σ : ℝ → ℝ) (hσ : ClassM σ) (hnp : ¬ IsAEPolynomial σ)
    (f : (Fin df → ℝ) → (Fin dm → ℝ) → ℝ)
    (hf : ContinuousOn (fun p => f p.1 p.2)
            (Set.Icc (0 : Fin df → ℝ) 1 ×ˢ Set.Icc (0 : Fin dm → ℝ) 1))
    (hmono : ∀ u ∈ Set.Icc (0 : Fin df → ℝ) 1,
        ∀ ⦃x y⦄, x ∈ Set.Icc (0 : Fin dm → ℝ) 1 → y ∈ Set.Icc (0 : Fin dm → ℝ) 1 →
          x ≤ y → f u x ≤ f u y)
    {ε : ℝ} (hε : 0 < ε) :
    ∃ P : PartMonoNet df dm, P.mono.IsMonotone ∧
      (∀ i, (fun u => P.emb u i) ∈ genSpanPi σ df) ∧
      ∀ u ∈ Set.Icc (0 : Fin df → ℝ) 1, ∀ x ∈ Set.Icc (0 : Fin dm → ℝ) 1,
        |P.toFun u x - f u x| ≤ ε
\end{verbatim}
\begin{proof}\leanok
  \uses{lem:leshno-bridge, thm:mono-approx}
  Fix $\varepsilon$. By uniform continuity of $f$ on the compact product, choose a fine tent
  partition of unity $\{\psi_i\}$ over a grid of the $u$-cube with support radius small enough that
  moving $u$ within a cell shifts $f$ by less than $\varepsilon/3$ (uniformly in $x$). With
  $g_i(x)=f(u_i,x)+C$ shifted nonnegative, the target $F(z,x)=\big(\sum_i z_i\,g_i(x)\big)-C$ is
  jointly continuous and coordinatewise monotone, so \cref{thm:mono-approx} gives a monotone
  network within $\varepsilon/3$ of $F$ on the cube. Each $\psi_i$ is approximated within a chosen
  $\eta$ by an unconstrained embedding coordinate via \cref{lem:leshno-bridge}, clamped into the
  cube; since $F$ is Lipschitz in $z$ this contributes another $\varepsilon/3$, and the partition
  collapse (the convex combination $\sum_i\psi_i(u)f(u_i,x)$) is within $\varepsilon/3$ of
  $f(u,x)$. The triangle inequality over the three terms yields $\varepsilon$. Soundness
  (\cref{thm:runje-sound}) makes the resulting network monotone in $x$.
\end{proof}

\subsection{Deep-core variant}

\begin{theorem}[Deep-core partial-monotone UAP]\label{thm:deep-part-uap}
  \lean{UniversalApproximation.Runje.deep_partial_monotone_approximation}\leanok
  \uses{thm:deep-mono-uap, def:partmononet}
  Swapping the shallow \texttt{MonoNet} core of \texttt{PartMonoNet} for a deep-residual
  \texttt{DeepMonoNet} of any depth (a \texttt{DeepPartMonoNet}) retains both soundness (monotone in
  the monotone block $x$, for every fixed $u$) and partial-monotone UAP: the witness embeds the
  shallow core of \cref{thm:runje-uap} via the single-block embedding of
  \cref{thm:deep-mono-uap}, sharing the same embedding, so the denotation and the approximation
  bound transfer verbatim.
\end{theorem}
